\documentclass[11pt]{article}
\usepackage[letterpaper,margin=0.85in]{geometry}
\usepackage{amsmath,amssymb,amsthm}
\usepackage{xcolor}
\usepackage[colorlinks=true,linkcolor=blue,urlcolor=blue]{hyperref}
\setlength{\parindent}{0pt}
\setlength{\parskip}{6pt}
\setlength{\emergencystretch}{2em}

\newtheorem{prop}{Proposition}
\newtheorem{cor}{Corollary}

\title{G1a.2: The Second-Invariant Scale Theorem\\
\large Conserved twist flux closes the two-parameter runaway}
\author{}
\date{2026-08-07 07:45 CST (UTC+8)}

\begin{document}
\maketitle

\section{Outcome}

The first two-parameter G1a model failed because frozen axial flux penalized small
radius but did not penalize the simultaneous deformation $r\to\infty$, $L\to0$.
Adding a second conserved magnetic structure with energy $K/L$ removes that runaway.

\begin{center}
\fbox{\parbox{0.9\textwidth}{
\textbf{Result.} If axial tube flux is frozen ($B>0$), a twist/poloidal-flux
invariant supplies $K/L$ with $K>0$, and either surface or core energy grows with
radius ($D>0$ or $F>0$), then the common-radius/common-length scale model has exactly
one finite equilibrium. It is the strict global energy minimum, and its physical
$2\times2$ Hessian is positive definite.
}}
\end{center}

This licenses finite scale in the idealized model. It does not yet prove a physical
theta-network equilibrium or branch stability.

\section{Why fixed twist flux gives \texorpdfstring{$K/L$}{K/L}}

For a tube of radius $r$ and length $L$, let $\Psi_\theta$ denote flux associated
with an azimuthal or twist field through a longitudinal surface. At the scaling
level,
\[
\Psi_\theta\sim B_\theta rL,
\qquad
E_\theta\sim\frac{B_\theta^2}{2\mu_0}r^2L
\sim\frac{\Psi_\theta^2}{2\mu_0L}
\equiv\frac{K}{L}.
\]
For two independent graph cycles, collect the conserved fluxes into
$\boldsymbol\Psi\in\mathbb R^2$. If the inductance operator scales as
$\mathsf M(L)=L\mathsf M_0$ with $\mathsf M_0$ symmetric positive definite, then
\[
E_{\rm cycle}
=\frac12\boldsymbol\Psi^{\mathsf T}\mathsf M(L)^{-1}\boldsymbol\Psi
=\frac{1}{2L}
\boldsymbol\Psi^{\mathsf T}\mathsf M_0^{-1}\boldsymbol\Psi
\equiv\frac{K}{L},
\qquad K>0.
\]
Thus mutual coupling between the theta graph's two cycles changes the positive
coefficient $K$ but not the scale exponent in this idealization.

The assumption that $\mathsf M_0$ is independent of the aspect ratio $L/r$ is a
declared limitation. Aspect-dependent and junction inductances are the next
refinement.

\section{Energy}

Retain the terms of the failed G1a model and add the second invariant:
\begin{equation}
E(r,L)=
AL+DrL+Fr^2L+B\frac{L}{r^2}
+\frac{K}{L}
+C(r^2L)^{-2/3}.
\label{eq:E}
\end{equation}
Here $r,L>0$; $A,C\ge0$; $B,K>0$; and at least one of $D,F$ is positive. The plasma
term uses fixed content and $\gamma=5/3$.

\section{Strict convexity in logarithmic scales}

Set $x=\log r$ and $y=\log L$. Equation~\eqref{eq:E} becomes
\[
\mathcal E(x,y)=
Ae^y+De^{x+y}+Fe^{2x+y}+Be^{-2x+y}
+Ke^{-y}+Ce^{-4x/3-2y/3}.
\]
Every term has the form $c_i e^{\mathbf v_i\cdot\mathbf z}$ with $c_i>0$ and
$\mathbf z=(x,y)$. Therefore
\[
\nabla^2_{\mathbf z}\mathcal E
=\sum_i c_ie^{\mathbf v_i\cdot\mathbf z}
\mathbf v_i\mathbf v_i^{\mathsf T}.
\]
The exponent vectors for the axial-flux and twist-flux terms,
$(-2,1)$ and $(0,-1)$, already span $\mathbb R^2$. Their positive outer-product
sum makes the logarithmic Hessian positive definite everywhere. Hence
$\mathcal E$ is strictly convex and can have at most one critical point.

\section{Existence and global minimality}

\begin{prop}[Second-invariant scale theorem]
Let $B,K>0$, $A,C\ge0$, and suppose $D>0$ or $F>0$. Then the energy in
Eq.~\eqref{eq:E} is coercive in $(x,y)=(\log r,\log L)$ and strictly convex.
Consequently it possesses exactly one critical point on $(r,L)\in(0,\infty)^2$,
and that point is the strict global minimum.
\end{prop}

\begin{proof}
Strict convexity was proved above. For coercivity, use the three exponent vectors
\[
\mathbf v_B=(-2,1),\qquad
\mathbf v_K=(0,-1),\qquad
\mathbf v_Q=(q,1),
\]
where $q=1$ if $D>0$ and $q=2$ if $F>0$. The origin is a strict convex combination
of these vectors: taking weights proportional to
\[
1,\qquad 1+\frac{2}{q},\qquad \frac{2}{q}
\]
gives zero, and all weights are positive. Thus the origin lies in the interior of
their convex hull. Along every unbounded direction in $(x,y)$ at least one of the
three corresponding exponentials diverges. Therefore $\mathcal E\to\infty$ as
$\|(x,y)\|\to\infty$.

A continuous coercive function attains a minimum. Strict convexity makes it unique.
Because the domain in logarithmic variables is all of $\mathbb R^2$, the minimizer
is a critical point.
\end{proof}

\begin{cor}
The Hessian with respect to physical scales $(r,L)$ is positive definite at the
critical point.
\end{cor}

\begin{proof}
At a critical point, the logarithmic and physical Hessians obey
\[
H_{\log}=
\begin{pmatrix}r&0\\0&L\end{pmatrix}
H_{\rm phys}
\begin{pmatrix}r&0\\0&L\end{pmatrix}.
\]
This is a congruence by a nonsingular positive diagonal matrix. Since $H_{\log}$ is
positive definite, so is $H_{\rm phys}$.
\end{proof}

\section{Closed-form control case}

When $C=0$, minimization over length gives
\[
L_*(r)=\sqrt{\frac{K}{g(r)}},
\qquad
E_*(r)=2\sqrt{Kg(r)},
\qquad
g(r)=A+Dr+Fr^2+\frac{B}{r^2}.
\]
The unique radius is the positive root of
\[
g'(r)=D+2Fr-\frac{2B}{r^3}=0,
\qquad
2Fr^4+Dr^3-2B=0.
\]
The left side is strictly increasing from $-2B$ to $+\infty$, proving existence and
uniqueness directly.

\section{Interpretation}

The two conserved structures do complementary work:
\[
\frac{BL}{r^2}\quad\text{punishes radial collapse},\qquad
\frac{K}{L}\quad\text{punishes longitudinal collapse}.
\]
Surface or core energy prevents unlimited radial expansion. Together they close the
two-dimensional scale plane.

This makes the second invariant a mathematical requirement of the current scale
model, not decorative reactor complexity. The theta graph's two independent cycles
are therefore potentially functional: their flux content can help select the
network scale.

\section{What this does not prove}

\begin{itemize}
\item It does not derive an order-parameter medium realizing the two fluxes.
\item It assumes common branch radius and a fixed dimensionless graph shape.
\item It compresses the cycle inductance matrix into a scale-independent positive
coefficient $K$; realistic aspect dependence may shift or remove the minimum.
\item It does not test junction algebra, zipping, kink/sausage modes, reconnection,
transport, or fusion benefit.
\item It is an existence theorem for an effective energy, not a reactor design.
\end{itemize}

\section{Verification}

Exact derivatives, exponent-vector convexity, the $C=0$ control solution, a
$C>0$ numerical example, and the physical Hessian are checked in
\path{../code/g1a_second_invariant_checks.py}.



\end{document}
